\documentclass[11pt]{article}
\usepackage[T1]{fontenc}
\usepackage{lmodern}
\usepackage{amsmath,amssymb,amsthm}
\usepackage[margin=1in]{geometry}
\usepackage{booktabs}
\usepackage[hidelinks]{hyperref}
\hypersetup{pdftitle={Negative Eliahou numbers with consecutive first generators},pdfauthor={Zhengyi Liu}}
\setlength{\emergencystretch}{2em}
\newtheorem{theorem}{Theorem}
\newtheorem{lemma}[theorem]{Lemma}
\newcommand{\N}{\mathbb N}
\newcommand{\card}[1]{\lvert #1\rvert}
\title{Negative Eliahou numbers with consecutive first generators}
\author{Zhengyi Liu}
\date{November 26, 2026}
\begin{document}
\maketitle
\begin{abstract}
We construct an infinite family of numerical semigroups whose first two
minimal generators are consecutive and whose Eliahou number is $-8$.
This disproves Conjecture 2.7 of Bacher, which asserts nonnegativity when
the second generator is at most the multiplicity plus six.
The proof uses a positional encoding to make short factorizations unique,
followed by a closed combinatorial count. No exhaustive search of
semigroups is needed.
\end{abstract}

\section{Statement}
For a numerical semigroup $S\subseteq\N$, let $m$ be its multiplicity,
$c$ its actual conductor, and $P$ its set of minimal generators, also
called primitive elements. Put
\[
 L=S\cap[0,c),\quad k=\card L,\quad \ell=\card{P\cap L},\quad
 q=\lceil c/m\rceil,\quad \rho=qm-c.
\]
Writing $s$ for the number of decomposable elements in $[c,c+m)$,
the Eliahou number is
\begin{equation}\label{eq:E}
 E(S)=k\ell-qs+\rho.
\end{equation}
Equivalently, if $r$ counts the primitive elements in that interval,
then $E(S)=k\ell+qr-c$, since $r+s=m$.

In the notation $m=\gamma_0<\gamma_1<\cdots$ for the left generators,
Bacher's Conjecture 2.7 states that $\gamma_1\le m+6$ implies
$E(S)\ge0$~\cite{Bacher}. The following family contradicts that assertion
even with $\gamma_1=m+1$.

\begin{theorem}\label{thm:main}
Let $M\ge10^6$ be an integer and define
\[
 G_M=\{4M,4M+1\}\cup\{5M+10^i:1\le i\le5\},
 \qquad S_M=\langle G_M\rangle\cup[32M,\infty).
\]
Then $S_M$ is a numerical semigroup with actual conductor $c=32M$,
multiplicity $m=4M$, and first two primitive elements $m,m+1$.
Moreover,
\[
 k=2024,\qquad \ell=7,\qquad s=1772,\qquad E(S_M)=-8.
\]
\end{theorem}

For instance, $M=10^6$ gives the explicit presentation
\begin{align*}
 S=\langle&4\,000\,000,\ 4\,000\,001,\ 5\,000\,010,\ 5\,000\,100,\\
           &5\,001\,000,\ 5\,010\,000,\ 5\,100\,000\rangle_{32\,000\,000}.
\end{align*}
Here the subscript means adjoining every integer starting at that value;
the proof below establishes that it is the actual conductor.

\section{Short factorizations and the actual conductor}
Write $T=\langle G_M\rangle$. A factorization in $T$ has the form
\begin{align}
 N&=4Mx_0+(4M+1)x_1+\sum_{i=1}^5(5M+10^i)v_i\notag\\
  &=M Q+\delta,\label{eq:encoding}\\
 Q&=4(x_0+x_1)+5t,\quad t=\sum_{i=1}^5v_i,
 \qquad \delta=x_1+\sum_{i=1}^5 10^i v_i,\notag
\end{align}
where all coefficients are nonnegative integers.

\begin{lemma}\label{lem:encoding}
The map in \eqref{eq:encoding} gives a bijection from coefficient
vectors satisfying $Q\le35$ to $T\cap[0,36M)$.
For such a vector, $N<32M$ if and only if $Q\le31$.
\end{lemma}
\begin{proof}
If $N<36M$, every summand is at least $4M$, so the total coefficient
$x_0+x_1+t$ is at most eight. Consequently
\[
 0\le\delta\le10^5(x_0+x_1+t)\le800000<M.
\]
Thus $Q=\lfloor N/M\rfloor\le35$.
Conversely, $Q\le35$ implies $4(x_0+x_1+t)\le35$, hence the same
coefficient and offset bounds, which give $N<36M$.
In either direction, $N<32M$ is equivalent to $Q\le31$.

For injectivity, $N$ uniquely determines its quotient $Q$ and remainder
$\delta$ upon division by $M$. Each of $x_1,v_1,\ldots,v_5$ is at most
eight. They are therefore precisely the six decimal digits of $\delta$,
with leading zeroes permitted. The quotient $Q$ then determines $x_0$.
\end{proof}

The set $S_M$ contains zero, is additively closed, and contains the entire
tail $[32M,\infty)$; it is therefore a numerical semigroup.
Its least positive element is $4M$.
If $32M-1$ belonged to $T$, Lemma~\ref{lem:encoding} would force its
remainder modulo $M$ to be at most $800000$. Its remainder is instead
$M-1\ge999999$. Thus $32M-1\notin S_M$, proving that $32M$ is the
actual conductor, not merely an upper bound for it.

All seven members of $G_M$ lie in $[4M,8M)$ and are distinct, so they
are primitive. Every other positive element of $T$ has a factorization
using at least two generators. Below $32M$ it is therefore decomposable
in $S_M$. Hence these seven elements are exactly the left primitives,
and the first two are $4M$ and $4M+1$.

\section{The two counts}
Set $r=x_0+x_1$. For fixed $t$, the number of choices of
$(v_1,\ldots,v_5)$ is $\binom{t+4}{4}$, and for fixed $r$ the number
of choices of $(x_0,x_1)$ is $r+1$.
Lemma~\ref{lem:encoding} now gives
\begin{equation}\label{eq:left}
 k=\sum_{t=0}^{6}\binom{t+4}{4}
       \sum_{r=0}^{\lfloor(31-5t)/4\rfloor}(r+1)
   =\sum_{t=0}^{6}\binom{t+4}{4}
       \binom{9-\lfloor5t/4\rfloor}{2}.
\end{equation}

The decomposable elements in $[32M,36M)$ are exactly
$T\cap[32M,36M)$. Indeed, a decomposition there cannot use an element
of the adjoined tail: its sum with any positive element would be at
least $32M+4M$. Both summands must therefore belong to $T$.
Conversely, an element of $T$ in this interval is larger than every
member of $G_M$, so any factorization of it uses at least two
generators and supplies a decomposition.

For each $0\le t\le7$, the interval
$32\le4r+5t\le35$ contains exactly one allowed value of $r$, namely
$r=8-\lfloor5t/4\rfloor$. Thus
\begin{equation}\label{eq:right}
 s=\sum_{t=0}^{7}\binom{t+4}{4}
                 \bigl(9-\lfloor5t/4\rfloor\bigr).
\end{equation}
The short sums \eqref{eq:left}--\eqref{eq:right} evaluate as follows.
These are counts of all coefficient vectors by degree, not a search
over bounded semigroups.
\begin{center}
\begin{tabular}{rrrr}
\toprule
$t$ & $\binom{t+4}{4}$ & contribution to $k$ & contribution to $s$\\
\midrule
0 & 1 & 36 & 9\\
1 & 5 & 140 & 40\\
2 & 15 & 315 & 105\\
3 & 35 & 525 & 210\\
4 & 70 & 420 & 280\\
5 & 126 & 378 & 378\\
6 & 210 & 210 & 420\\
7 & 330 & 0 & 330\\
\midrule
Total & & 2024 & 1772\\
\bottomrule
\end{tabular}
\end{center}
Since $c=8m$, we have $q=8$ and $\rho=0$. Substitution into
\eqref{eq:E} completes the proof of Theorem~\ref{thm:main}:
\[
 E(S_M)=7\cdot2024-8\cdot1772=-8.
\]

\section{Scope and reproducibility}
Theorem~\ref{thm:main} disproves the small-gap implication of
Conjecture 2.7 in version 9 of~\cite{Bacher}. It does not disprove
Wilf's conjecture. Indeed, the number of right primitives is
$4M-1772$, so the total embedding dimension is $4M-1765$ and
\[
 W(S_M)=2024(4M-1765)-32M=8064M-3572360>0.
\]
No minimality claim is made. The large decimal separation makes the
proof transparent. Supplementary exact calculations accompany the
uniform symbolic proof above.

\paragraph{Formal verification.}
The companion Lean 4 development proves the full family from intrinsic
semigroup definitions and refutes the universally quantified conjecture,
using only Lean's three standard axioms. The
\href{https://github.com/Onetkn/eliahou-small-gap-counterexample}{source and verification records}
are available online.

\paragraph{AI assistance.}
OpenAI Codex assisted in the construction, proof development, computational
checks, formalization, and manuscript preparation.

\begin{thebibliography}{9}
\bibitem{Bacher}
A. Bacher, \emph{Families of Eliahou semigroups linked to Farey intervals},
arXiv:2604.25051v9, 14 October 2026.
\url{https://arxiv.org/abs/2604.25051v9}.
\end{thebibliography}
\end{document}
